\chapter{Market making sur flux dans une banque}\label{ch:s2:flow-market-making-at-a-bank}

Le desk de flux d'une banque gagne de l'argent en valorisant ses clients, et non en prédisant les cours. Il apprend quelles transactions de ses clients vont se retourner contre lui et leur fait payer
en conséquence, et il garde le risque de la transaction d'un client jusqu'à ce qu'une transaction d'un autre client la compense. Butz et Oomen ont estimé qu'un grand dealer représentatif
met plusieurs minutes en moyenne à internaliser la transaction d'un client sur les devises les plus liquides, et des dizaines de minutes sur les marchés émergents.
Sur le desk synthétique de ce chapitre, valoriser chaque client à partir de son propre historique fait passer le \omterm{def:s2:flow-market-making-at-a-bank:franchise}{P\&L de franchise} de 6.63~millions de dollars, obtenu avec le meilleur
de quatre prix uniformes, à 7.60~millions de dollars. Les clients non informés paient 1.39~points de base et les clients informés cessent de trader. Un petit \omterm{def:s2:flow-market-making-at-a-bank:skew}{décalage de prix guidé par les axes} fait passer la part du risque
internalisée de 91~\% à 97~\%. Le code est \texttt{firm.flowmm}.

\section{La franchise et son P\&L}

\begin{definition}[P\&L de franchise]\label{def:s2:flow-market-making-at-a-bank:franchise}
Le \emph{P\&L de franchise}\index{P\&L de franchise} d'un desk de flux est son profit tiré du service rendu aux clients, décomposé en capture du spread (le prix que paient les clients
paient par rapport au mid), la valorisation de l'inventaire que laissent leurs transactions (où se manifeste la sélection adverse), et le coût de la couverture de ce
qu'il n'internalise pas.
\end{definition}

\texttt{firm.flowmm} envoie à un desk 200\,000 demandes de cotation (RFQ) de 90 clients, chacune portant sur une unité de 1~million de dollars, de sorte qu'un point de base vaut 100~dollars. Soixante
clients sont non informés~: après leurs transactions, le mid évolue au hasard. Vingt sont modérément informés~: le mid dérive de 0.8~points de base dans leur sens
sur les dix demandes suivantes. Dix sont informés, avec une dérive de 2.5~points de base. La dérive se produit quel que soit celui qui exécute la transaction, de sorte que la trajectoire du mid est la même
quel que soit le prix que cote le desk, et chaque politique de valorisation est comparée sur les mêmes prix. Chaque client dispose aussi de sa propre cotation concurrente, le demi-spread
qu'il peut obtenir ailleurs~: 1.0 à 2.5~points de base pour les clients non informés, dont certains sont captifs, 0.8 à 1.5 pour les modérés et 0.4 à 0.8 pour les
informés, qui font jouer la concurrence. Un client accepte le prix du desk avec une probabilité qui décroît régulièrement à mesure que ce prix dépasse sa cotation concurrente
(voir le \cref{lst:s2:flow-market-making-at-a-bank:desk}).

Le desk cote tout le monde à 1.2~points de base pour les 100\,000 premières demandes et apprend de celles-ci. Le mark-out, soit le déplacement du mid au cours des 20
demandes suivantes, contre le côté du desk, sépare les trois types (voir la \cref{fig:s2:flow-market-making-at-a-bank:markout})~: $-0.01$, 0.84 et 2.41~points de base en moyenne.
Le taux de touche (hit ratio), c'est-à-dire la part des demandes d'un client que le desk remporte, révèle sa cotation concurrente~: 80~\%, 43~\% et 8~\%. Inférées à partir du
taux de touche de chaque client, les cotations concurrentes présentent une corrélation de 0.998 avec les véritables.

\begin{omfigure}
\begin{tikzpicture}
\begin{axis}[width=11cm,height=5.2cm,xlabel={\small demandes après la transaction},ylabel={\small mark-out (pb, en faveur du client)},
  tick label style={font=\footnotesize},xmin=0,xmax=40,ymin=-0.5,ymax=3,grid=major,grid style={gray!25},table/col sep=comma,
  legend style={font=\scriptsize,at={(0.5,-0.3)},anchor=north},legend columns=3]
\addplot[black,thick] table[x=horizon,y=uninformed]{figdata/strategies-2/24-flow-market-making-at-a-bank/markout.csv};
\addplot[omDef,thick,dashed] table[x=horizon,y=moderate]{figdata/strategies-2/24-flow-market-making-at-a-bank/markout.csv};
\addplot[omThm,thick] table[x=horizon,y=informed]{figdata/strategies-2/24-flow-market-making-at-a-bank/markout.csv};
\legend{non informés,modérément informés,informés}
\end{axis}
\end{tikzpicture}

\omcaption{Courbes de mark-out des transactions exécutées du desk synthétique, par type de client~: le déplacement du mid après la transaction, en faveur du client, exprimé en points
de base, en fonction du nombre de demandes écoulées depuis. Données~: \texttt{s2\_flowmm.markout\_curve}.}\label{fig:s2:flow-market-making-at-a-bank:markout}
\end{omfigure}

\section{Segmentation des clients (tiering)}

\begin{definition}[Segmentation des clients (tiering)]\label{def:s2:flow-market-making-at-a-bank:tiering}
La \emph{segmentation des clients}\index{segmentation des clients} (tiering) consiste à valoriser chaque client, ou groupe de clients, selon ce que ses transactions coûtent au desk (son mark-out) et selon
sa sensibilité au prix (son taux de touche face aux cotations du desk), plutôt que de coter tout le monde de la même façon.
\end{definition}

Pour chaque client, le desk choisit le demi-spread qui maximise son profit espéré par demande, soit le demi-spread diminué du mark-out du client, multiplié par la
probabilité de remporter la demande à ce prix (voir le \cref{lst:s2:flow-market-making-at-a-bank:price}). Les clients captifs non informés paient davantage, les clients sensibles au
prix paient moins, et les clients informés sont écartés du marché. Sur les 100\,000 demandes suivantes (P\&L en millions de dollars)~:

\begin{center}
\small
\setlength{\tabcolsep}{4pt}%
\begin{tabular}{@{}p{5.0cm}ccccc@{}}
\toprule
politique & transactions & capture & mark-out & couverture & \omterm{def:s2:flow-market-making-at-a-bank:franchise}{P\&L de franchise}\\
\midrule
prix uniforme 1.0~pb & 74\,622 & 7.46 & $-1.21$ & 0.35 & 5.90\\
prix uniforme 1.2~pb & 64\,307 & 7.72 & $-0.84$ & 0.31 & 6.57\\
prix uniforme 1.4~pb & 53\,310 & 7.46 & $-0.58$ & 0.25 & 6.63\\
prix uniforme 1.6~pb & 42\,399 & 6.78 & $-0.36$ & 0.20 & 6.23\\
prix par client & 62\,147 & 8.53 & $-0.65$ & 0.28 & 7.60\\
prix par client, décalage guidé par les axes 0.05 & 61\,519 & 8.35 & $-0.62$ & 0.10 & 7.63\\
\bottomrule
\end{tabular}
\end{center}

Aucun prix uniforme ne fait aussi bien. Un prix bas remporte les transactions des clients informés et y perd de l'argent~; un prix élevé fait perdre les affaires des clients captifs, en même temps que
les clients informés. Par type, du prix uniforme à 1.2~points de base aux prix par client~:

\begin{center}
\small
\begin{tabular}{@{}lcccc@{}}
\toprule
par transaction (pb) & transactions & capture & mark-out & net\\
\midrule
non informés, uniforme & 53\,500 & 1.20 & 0.00 & 1.20\\
non informés, par client & 53\,600 & 1.39 & 0.00 & 1.40\\
modérés, uniforme & 9\,869 & 1.20 & 0.83 & 0.37\\
modérés, par client & 8\,537 & 1.25 & 0.82 & 0.44\\
informés, uniforme & 938 & 1.20 & 2.38 & $-1.18$\\
informés, par client & 10 & & &\\
\bottomrule
\end{tabular}
\end{center}

Oomen a montré que les signatures de prix, soit la trajectoire moyenne du prix autour des exécutions, permettent de distinguer les fournisseurs de liquidité qui internalisent de ceux qui externalisent~;
les mêmes courbes, calculées par client, sont celles sur lesquelles un desk fonde sa segmentation. La segmentation a un pendant~: la qualité d'exécution d'un client dépend de la manière dont ses dealers
gèrent le risque. Oomen a étudié le last look, soit l'option du dealer de rejeter une demande de transaction après avoir vu où le prix s'est porté, et a constaté que la conception
de cette option détermine le risque d'exécution du client mais ne modifie pas forcément ses coûts effectifs.

\section{Axes et décalage de prix}

\begin{definition}[Décalage de prix guidé par les axes]\label{def:s2:flow-market-making-at-a-bank:skew}
Un \emph{décalage de prix guidé par les axes}\index{décalage de prix guidé par les axes} est un déplacement des deux côtés de la cotation d'un dealer en fonction de son inventaire ou de ses axes, de sorte que les clients
se voient proposer de meilleurs prix sur les transactions qui réduisent la position du desk et de moins bons sur celles qui l'augmentent.
\end{definition}

Le desk décale ses cotations de 0.05~points de base par unité d'inventaire~: long de cinq unités, il vend 0.25~points de base moins cher et achète 0.25~points de base
moins cher aussi. Les clients échangent alors davantage dans le sens qui réduit la position, et le desk se couvre moins. Le décalage coûte en capture et fait économiser de la
couverture et du risque (voir la \cref{fig:s2:flow-market-making-at-a-bank:skew})~: à 0.05, le \omterm{def:s2:flow-market-making-at-a-bank:franchise}{P\&L de franchise} est légèrement plus élevé, à 7.63~millions de dollars, et l'inventaire moyen
passe de 2.75 à 2.18 unités~; à 0.2, le desk ne se couvre jamais, détient en moyenne 1.30 unité et gagne 7.08~millions de dollars.

\begin{omfigure}
\begin{tikzpicture}
\begin{axis}[width=11cm,height=5.2cm,axis y line*=right,axis x line=none,ylabel={\small part internalisée (\%)},
  tick label style={font=\footnotesize},xmin=0,xmax=0.2,ymin=88,ymax=101,table/col sep=comma]
\addplot[omDef,thick,dashed,mark=square*,mark size=1pt] table[x=skew,y=internalised]{figdata/strategies-2/24-flow-market-making-at-a-bank/skew.csv};
\label{plot:s2:flowint}
\end{axis}
\begin{axis}[width=11cm,height=5.2cm,axis y line*=left,xlabel={\small décalage guidé par les axes (pb par unité d'inventaire)},
  ylabel={\small P\&L de franchise (milliers de pb)},tick label style={font=\footnotesize},xmin=0,xmax=0.2,ymin=70,ymax=77,grid=major,
  grid style={gray!25},table/col sep=comma,scaled x ticks=false,xticklabel style={/pgf/number format/fixed},
  legend style={font=\scriptsize,at={(0.5,-0.3)},anchor=north},legend columns=2]
\addplot[omThm,thick,mark=*,mark size=1pt] table[x=skew,y=net]{figdata/strategies-2/24-flow-market-making-at-a-bank/skew.csv};
\addlegendentry{P\&L de franchise (gauche)}
\addlegendimage{/pgfplots/refstyle=plot:s2:flowint}\addlegendentry{part internalisée (droite)}
\end{axis}
\end{tikzpicture}

\omcaption{Le desk synthétique tarifé par client, seconde moitié~: \omterm{def:s2:flow-market-making-at-a-bank:franchise}{P\&L de franchise} (en milliers de points de base de notionnel unitaire, 100~\$ chacun) et la part
du volume client internalisé, en fonction du décalage guidé par les axes. Données~: \texttt{s2\_flowmm.policy}.}\label{fig:s2:flow-market-making-at-a-bank:skew}
\end{omfigure}

\section{Internalisation}

Un desk internalise la transaction d'un client lorsqu'il conserve le risque jusqu'à ce qu'une autre transaction de client la compense, et il externalise le risque lorsqu'il couvre sur le
marché. Le desk synthétique détient jusqu'à cinq unités et couvre au-delà de cette limite au demi-spread de marché de 0.5 point de base. Tarifé par client sans
skew, il internalise 90.8~\% de son volume~; avec le skew de 0.05, 96.8~\%. Butz et Oomen ont modélisé l'internalisation comme une file d'attente et ont constaté que ses coûts
sont plus faibles pour les dealers prêts à détenir davantage de risque et pour ceux dont les clients sont plus sensibles au prix. Ils concluent que les coûts d'un client dépendent
de la gestion du risque du dealer autant que de son propre trading, et que les clients doivent distinguer les externalisateurs des internalisateurs passifs et agressifs.

\section{Fiches de stratégie}

\begin{strategyfile}{Tarification différenciée par mark-out client}\label{strat:s2:flow-market-making-at-a-bank:tier}
\sfield{Qui paie le desk, et pourquoi} Des clients qui valorisent la relation, la commodité ou le crédit plutôt que le dernier point de base.
\sfield{Instruments et plateformes} Demandes de cotation (RFQ) et cours diffusés aux clients, par plateforme ou par téléphone.
\sfield{Signal} Courbe de mark-out et taux de touche (hit ratio) de chaque client.
\sfield{Dimensionnement et exécution} Demi-spread par client qui maximise le profit espéré par demande.
\sfield{Coûts} Sélection adverse~; affaires perdues à cause d'un prix trop élevé.
\sfield{Comment elle meurt} Clients qui changent de comportement~; concurrents qui segmentent mieux.
\sfield{Horizon, capacité, infrastructure} Continu~; analytique client et moteurs de valorisation.
\sfield{Backtest honnête} Mark-outs sur les seules transactions gagnées~; la courbe de gain estimée à partir de cotations variées, non supposée.
\sfield{Sources} Oomen (2019)~; ce chapitre~: 7.60~millions de dollars contre 6.63~millions au meilleur prix plat.
\end{strategyfile}

\begin{strategyfile}{Décalage guidé par les axes}\label{strat:s2:flow-market-making-at-a-bank:skew}
\sfield{Qui paie le desk, et pourquoi} Les clients situés du côté opposé de la position du desk, qui obtiennent un meilleur prix.
\sfield{Instruments et plateformes} Les cotations du desk~; les axes affichés aux clients.
\sfield{Signal} Inventaire et axes.
\sfield{Dimensionnement et exécution} Un skew par unité d'inventaire, à mettre en balance avec la capture perdue et la couverture économisée.
\sfield{Coûts} Capture abandonnée.
\sfield{Comment elle meurt} Clients qui lisent le skew et négocient contre lui.
\sfield{Horizon, capacité, infrastructure} Continu.
\sfield{Backtest honnête} Réponse des clients au skew estimée, non supposée.
\sfield{Sources} Ce chapitre~: internalisation passant de 90.8~\% à 96.8~\% à 0.05.
\end{strategyfile}

\begin{strategyfile}{Internalisation avant couverture}\label{strat:s2:flow-market-making-at-a-bank:internalise}
\sfield{Qui paie le desk, et pourquoi} La fourchette du marché non versée sur les flux de sens opposé.
\sfield{Instruments et plateformes} L'inventaire propre du desk~; le marché pour couvrir le reliquat.
\sfield{Signal} Flux de sens opposé attendu, à comparer au risque d'attendre.
\sfield{Dimensionnement et exécution} Une limite de risque au-delà de laquelle le desk couvre.
\sfield{Coûts} Risque d'inventaire~; information divulguée par la couverture.
\sfield{Comment elle meurt} Marchés à sens unique.
\sfield{Horizon, capacité, infrastructure} Quelques minutes à quelques heures~; un central risk book (CRB).
\sfield{Backtest honnête} Valorisation de l'inventaire sur l'horizon d'internalisation.
\sfield{Sources} Butz et Oomen (2019).
\end{strategyfile}

\begin{strategyfile}{Gestion du taux de touche des RFQ}\label{strat:s2:flow-market-making-at-a-bank:winrate}
\sfield{Qui paie le desk, et pourquoi} Des clients dont l'activité vaut d'être gagnée à un certain prix.
\sfield{Instruments et plateformes} Plateformes de demande de cotation multi-dealers.
\sfield{Signal} Taux de touche par client, par taille et par moment, comparés aux fourchettes cotées.
\sfield{Dimensionnement et exécution} Taux de réussite cibles par segment de clientèle.
\sfield{Coûts} Malédiction du vainqueur sur les demandes informées.
\sfield{Comment elle meurt} Course au volume à mark-out négatif.
\sfield{Horizon, capacité, infrastructure} Continu~; analytique de plateforme.
\sfield{Backtest honnête} Les issues des demandes perdues ne sont pas observées~; ne pas les supposer.
\sfield{Sources} Aucun chiffre de performance vérifié~; taux de touche de ce chapitre~: 80~\%, 43~\% et 8~\%.
\end{strategyfile}

\section{Tutoriel~: tarifer le client}\label{tut:s2:flow-market-making-at-a-bank:tut}

\begin{tutorial}
\textbf{Objectif.} Simuler un desk face à des clients aux informations et à la concurrence différentes, apprendre leurs mark-outs et taux de touche, les tarifer, appliquer un skew
selon l'inventaire et décomposer le \omterm{def:s2:flow-market-making-at-a-bank:franchise}{P\&L de franchise}. \textbf{État final~:} les trois tableaux et les trois figures.
\begin{tutsteps}
\item \textbf{Le desk}.
\omcode{code/firm/flowmm/firm_flowmm.py}{63}{94}{Coter, exécuter, internaliser, couvrir au-delà de la limite, décomposer.}\label{lst:s2:flow-market-making-at-a-bank:desk}
\item \textbf{Apprendre et tarifer les clients}.
\omcode{code/firm/flowmm/firm_flowmm.py}{97}{130}{Mark-outs, taux de touche et prix qui maximise le profit espéré.}\label{lst:s2:flow-market-making-at-a-bank:price}
\item \textbf{Exécuter} \texttt{estimates()}, \texttt{policy()}, \texttt{by\_type()} et \texttt{fig\_flowmm.py}.
\end{tutsteps}
\textbf{Pour aller plus loin.} Faire changer les clients de type à mi-parcours~; estimer la courbe de gain à partir de cotations randomisées plutôt que supposer sa forme~; ajouter le
last look.
\end{tutorial}

\section{Construction~: market making sur flux}\label{bld:s2:flow-market-making-at-a-bank:flowmm}

\begin{build}
\bfield{Rôle} Un desk de RFQ avec mark-outs client, tarification par client, décalage guidé par les axes et internalisation, avec son P\&L décomposé.
\bfield{Interface} \texttt{FlowConfig(\dots)}, \texttt{simulate\_flow(cfg)}, \texttt{run\_desk(flow, cfg, half\_spread, skew, start, end)},
\texttt{client\_markouts}, \texttt{client\_hits}, \texttt{price\_clients}.
\bfield{Règles} Le chemin du prix milieu ne dépend pas du desk~; le choix du client est logistique en fonction du prix payé~; couverture seulement au-delà de la limite.
\bfield{Tests d'acceptation} \texttt{code/firm/flowmm/tests/}~: la décomposition s'additionne et des cotations plus larges sont moins exécutées~; les mark-outs classent les types~; la tarification
à la main~; inventaire dans la limite.
\bfield{Pour aller plus loin} Last look~; cours diffusés~; plusieurs produits à inventaires corrélés.
\end{build}

\begin{omsources}
\item M.~Butz et R.~Oomen, ``Internalisation by electronic FX spot dealers'', \emph{Quantitative Finance} 19(1), 2019.
\item R.~Oomen, ``Last look'', \emph{Quantitative Finance} 17(7), 2017.
\item R.~Oomen, ``Price signatures'', \emph{Quantitative Finance} 19(5), 2019.
\end{omsources}

\section{Exercices}

\begin{exercise}[$\star$]\label{exo:s2:flow-market-making-at-a-bank:1}
Un client paie 1.2 points de base sur 1~million de dollars, et ses transactions ont un mark-out de 0.83. Calculer le gain du desk par transaction, en dollars.
\end{exercise}

\begin{exercise}[$\star$]\label{exo:s2:flow-market-making-at-a-bank:2}
Le desk est long de 4 unités avec un skew de 0.05 points de base par unité et un demi-spread de 1.2. Calculer ce que paie un client acheteur, et ce que paie un client vendeur.
\end{exercise}

\begin{exercise}[$\star$]\label{exo:s2:flow-market-making-at-a-bank:3}
Le taux de touche d'un client à 1.2 points de base est de 80~\%, et le choix a une largeur de 0.25. Calculer la cotation concurrente que cela implique.
\end{exercise}

\begin{exercise}[$\star\star$]\label{exo:s2:flow-market-making-at-a-bank:4}
Pourquoi aucun prix plat ne rejoint-il la tarification par client~?
\end{exercise}

\begin{exercise}[$\star\star$]\label{exo:s2:flow-market-making-at-a-bank:5}
Pourquoi un skew faible fait-il monter le \omterm{def:s2:flow-market-making-at-a-bank:franchise}{P\&L de franchise} alors qu'un skew important le fait baisser~?
\end{exercise}

\begin{exercise}[$\star\star$]\label{exo:s2:flow-market-making-at-a-bank:6}
Pourquoi les coûts d'internalisation sont-ils plus faibles pour les dealers dont les clients sont plus sensibles au prix~?
\end{exercise}

\begin{exercise}[$\star\star\star$]\label{exo:s2:flow-market-making-at-a-bank:7}
\emph{Programmation.} Exécuter \texttt{policy} avec un demi-spread plat de 0.8 point de base. Que se passe-t-il pour les transactions, les valorisations et le \omterm{def:s2:flow-market-making-at-a-bank:franchise}{P\&L de franchise}~?
\end{exercise}

\begin{exercise}[$\star\star\star$]\label{exo:s2:flow-market-making-at-a-bank:8}
\emph{Trouver la faille.} «~Nos clients informés génèrent l'essentiel de la croissance de nos volumes~; nous devrions resserrer leurs spreads pour en gagner davantage.~»
\end{exercise}

\section{Problème~: tarification du client}

\begin{problem}[{Problème du week-end --- flow market making}]\label{pb:s2:flow-market-making-at-a-bank:1}
Le desk synthétique du chapitre et les données publiques.

\medskip\noindent\textbf{Partie I --- La franchise.}
\begin{enumerate}
\item Définir le \omterm{def:s2:flow-market-making-at-a-bank:franchise}{P\&L de franchise} et ses composantes.
\item Décrire les clients et expliquer pourquoi le chemin du mid ne dépend pas du desk.
\item Que montrent les mark-outs et les taux de touche (hit ratio) après l'historique à prix plat~?
\item Que montraient les signatures de prix d'Oomen~?
\end{enumerate}

\medskip\noindent\textbf{Partie II --- Segmentation (tiering).}
\begin{enumerate}[resume]
\item Définir la \omterm{def:s2:flow-market-making-at-a-bank:tiering}{segmentation des clients} (tiering).
\item Comment chaque client est-il tarifé~?
\item Donner le tableau des politiques.
\item Donner la capture du spread et le net par transaction, selon le type.
\end{enumerate}

\medskip\noindent\textbf{Partie III --- Skews et internalisation.}
\begin{enumerate}[resume]
\item Définir un \omterm{def:s2:flow-market-making-at-a-bank:skew}{décalage de prix guidé par les axes} (axe).
\item Quel effet le skew a-t-il sur le P\&L, l'inventaire et l'internalisation~?
\item Que trouvent Butz et Oomen~?
\item Qu'est-ce que le last look, et que trouve Oomen~?
\end{enumerate}

\medskip\noindent\textbf{Partie IV --- Le verdict.}
\begin{enumerate}[resume]
\item Énoncer le \emph{résultat nommé}~: la capture du spread du desk par segment et la part du risque internalisée.
\item Quelle hypothèse flatte le plus la tarification par client~?
\item Pourquoi un prix plat bas fait-il perdre de l'argent~?
\item Comment estimer honnêtement la courbe de gain d'un client~?
\item Que faire d'un client dont le mark-out augmente~?
\item Quelle fiche de stratégie porte le plus de risque d'inventaire~?
\item Quel lien existe-t-il entre ce chapitre et les stratégies de flux du chapitre~15~?
\item En une phrase~: que vend un desk de flux~?
\end{enumerate}
\end{problem}

\section{Questions d'entretien}

\begin{interviewq}[$\star$ \iqroles{trader}]\label{iq:s2:flow-market-making-at-a-bank:1}
Qu'est-ce qu'un mark-out, et comment l'utiliseriez-vous pour tarifer un client~?
\end{interviewq}

\begin{interviewq}[$\star\star$ \iqroles{chercheur}]\label{iq:s2:flow-market-making-at-a-bank:2}
Vous ne voyez que les demandes que vous avez remportées. Comment estimez-vous la sensibilité d'un client à votre prix~?
\end{interviewq}

\begin{interviewq}[$\star\star$ \iqroles{trader}]\label{iq:s2:flow-market-making-at-a-bank:3}
Vous êtes long d'une position importante et un client demande une cotation bilatérale. Que cotez-vous~?
\end{interviewq}

\begin{interviewq}[$\star\star$ \iqroles{risque}]\label{iq:s2:flow-market-making-at-a-bank:4}
Comment fixeriez-vous la limite d'inventaire d'un desk qui internalise~?
\end{interviewq}

\begin{interviewq}[$\star\star$ \iqroles{développeur}]\label{iq:s2:flow-market-making-at-a-bank:5}
Concevez le service qui calcule en temps réel les mark-outs des clients pour un moteur de valorisation.
\end{interviewq}

\begin{interviewq}[$\star\star\star$ \iqroles{chercheur}]\label{iq:s2:flow-market-making-at-a-bank:6}
Avec une probabilité de gain logistique $p(h) = 1/(1 + e^{(h - c)/w})$ et un mark-out $m$, montrez que le demi-spread maximisant le profit satisfait
$h - m = w / (1 - p(h))$.
\end{interviewq}
